% !TeX root = main.tex
\newpage
\section{Cathedral}

\subsection{Initial Lighting Setup}

The cathedral lighting scenario is a three-dimensional particle visualization
for light interacting with suspended dust.  The left side of the scene emits a
stream of photon particles.  The photon stream enters the domain at a downward
angle of approximately \(-15^\circ\), so the light crosses the space from left
to right while descending.

The bottom of the scene emits a second, slower stream of dust particles.  The
dust stream rises upward through the light path.  The intended visual target is
a visible illuminated volume where the downward photon stream intersects the
slow upward dust motion.

Both sources use the streaming generator's emitter-plane form.  A stream fixes
one inlet coordinate and spans a finite \(u,v\) patch on that plane.  The photon
beam is a rectangular patch on the left side of the scene, while the dust source
is a rectangular patch on the bottom of the scene.  The patch fields are a
generation contract: Python expands them into particle positions in the binary
data file, and the Vulkan runtime consumes the resulting particles rather than
the stream-shape fields themselves.

The initial cathedral wall model is deliberately minimal.  The scene uses a
single right-side wall so photons and dust have a visible boundary target
without adding the older reservoir helper walls.

\subsection{Light Particle Behavior}

Photon particles use the lumens material color behavior learned from the
earlier lumens prototype.  Lumens is a rendering rule, not a force rule:
photon particles still move according to their configured particle state and
dynamics, while visibility and brightness are chosen by the renderer from the
current particle state.

A photon becomes visibly lit only when its current position reaches the
configured eye or sensor aperture.  The renderer must not color a photon early
just because the future continuation of its ray would later pass through the
aperture.  This prevents wall rebounds or near misses from creating light
before the particle is actually in the visible region.

The lumens rule uses the particle's current position, current velocity angle,
the eye position, the eye look direction, the field of view, and the aperture
radius.  A particle contributes visible light only when:
\begin{itemize}
    \item its current position is inside the aperture radius;
    \item its velocity ray passes through the aperture rather than away from
          it;
    \item its arrival direction falls inside the eye field of view.
\end{itemize}

When \texttt{lumens\_fade\_in} is enabled, brightness ramps up as the photon
moves from the aperture boundary toward the eye center.  When it is disabled,
all valid photons inside the aperture use the angular brightness rule without
distance fade.  \texttt{lumens\_debug\_dim} may draw non-visible photons dimly
for diagnostics, but that is a viewer option and not part of the physical
lighting model.

\subsection{Photon Collision Rules}

Photons have a special interaction rule carried by material configuration.
Each material may declare a \texttt{particle\_type}; regular matter uses the
regular type, while the photon beam uses the photon type.  The generated
material data preserves that value for the shader path.

Photon particles do not collide with, rebound from, or exchange momentum with
other photon particles.  The shared particle-particle force loop skips
photon-photon pairs before contact evaluation.  A photon stream can therefore
pass through itself without creating photon-photon pressure, compression, or
scattering.

Photon speed is invariant.  The configured stream
\texttt{initial\_particle\_velocity} sets the photon speed by its vector
magnitude.  There is no separate material speed field.  When a photon collides
with dust or a wall, the collision may change the photon's direction, but the
outgoing photon velocity must be renormalized to the incoming photon speed.

Photons only collide with real material particles, such as the rising dust
particles in the cathedral scene.  Those interactions are the intended source
of visible lighting behavior: photons crossing the dust stream can mark,
illuminate, or otherwise interact with the real particles, while remaining
transparent to other photons.

Photon and dust particles are separate momentum systems.  A photon-dust
interaction redirects the photon only; the dust particle's velocity is
unchanged by that photon.  Dust momentum is governed by dust-dust and dust-wall
interactions, not by photon contacts.

Photon-wall interaction remains active.  The particle-pair skip applies to
photon-photon pairs and to regular-source/photon-target pairs.  Photon-source
and dust-target pairs remain active so the photon can reflect without applying
a reciprocal force to the dust.  Wall and boundary-marker paths are not
disabled by the photon rule.

Photons are also allowed to tunnel.  A photon may pass through intervening
space without requiring the ordinary particle rebound behavior used by real
matter particles.  Any photon collision or visibility rule must therefore be
explicit about which real-particle interactions are meaningful and must not
assume that photon motion is constrained by the same rebound rules as massive
particles.
